% Bagian: sets-functions-relations
% Bab: arithmetization
% Subbagian: pemeriksaan-rincian
%
\documentclass[../../../include/open-logic-section]{subfiles}


\begin{document}
	
\olfileid[id]{sfr}{arith}{check}
\olsection{Ring dan Lapangan Terurut}

Sepanjang bab ini, kita mengklaim bahwa definisi-definisi tertentu berperilaku
``sebagaimana semestinya''. Dalam lampiran teknis ini, kita akan menjabarkan
apa yang kita maksud, serta (menguraikan cara) menunjukkan bahwa
definisi-definisi itu memang berperilaku ``dengan tepat''.

Dalam \olref[int]{sec}, kita mendefinisikan penjumlahan dan perkalian pada
$\Int$. Kita ingin menunjukkan bahwa, sebagaimana didefinisikan, keduanya
memberi $\Int$ struktur yang ``kita inginkan''. Secara khusus, struktur yang
dimaksud adalah ring komutatif.

\begin{defn}
	Suatu \emph{ring komutatif} adalah himpunan $S$, yang dilengkapi elemen-elemen khusus $0$ dan $1$ serta operasi $+$ dan $\times$, yang memenuhi kedelapan rumus berikut:
	\begin{align*}
		\emph{Asosiativitas}&&a + (b+ c) & = (a + b) + c \\
		&& (a \times b) \times c & = a \times (b\times c)\\
		\emph{Komutativitas}&&a + b &= b+ a  \\
		&&  a \times b&= b\times a\\ 
		\emph{Elemen identitas}&&a + 0 &= a \\
		&& a \times 1 &= a\\
		\emph{Invers aditif}&&(\exists b\in S)0&=a + b\\
		\emph{Distributivitas}&&a \times (b+ c ) &= (a \times b) + (a \times c)
	\end{align*}
	Secara implisit, semua rumus ini terikat oleh kuantor universal yang dibatasi pada $S$. Perhatikan pula bahwa elemen $0$ dan~$1$ di sini tidak harus merupakan bilangan asli yang bernama sama.
\end{defn}

Jadi, untuk memeriksa bahwa bilangan bulat membentuk suatu ring komutatif,
kita hanya perlu memeriksa bahwa kedelapan kondisi tersebut terpenuhi. Tidak
ada kondisi yang {sulit} dibuktikan, tetapi pemeriksaannya agak panjang.
Sebagai contoh, berikut cara membuktikan \emph{Asosiativitas} untuk
penjumlahan:

\begin{proof} 
Tetapkan $i, j, k \in \Int$. Jadi terdapat $a_1, b_1, a_2, b_2, a_3, b_3 \in
\Nat$ sedemikian sehingga $i = \equivrep{a_1, b_1}{}$ dan $j =
\equivrep{a_2,b_2}{}$ dan $k = \equivrep{a_3, b_3}{}$. (Agar mudah dibaca,
kita menulis ``$\equivrep{x, y}{}$'' alih-alih
``$\equivrep{x, y}{\Intequiv}$''; kita akan melakukannya sepanjang
bagian ini.) Sekarang:
\begin{align*}
	i + (j + k) &= \equivrep{a_1, b_1}{}+(\equivrep{a_2, b_2}{} + \equivrep{a_3, b_3}{}) \\
	&= \equivrep{a_1,  b_1}{} + \equivrep{a_2+a_3, b_2+b_3}{}\\
	&= \equivrep{a_1 + (a_2 + a_3), b_1 + (b_2 + b_3)}{}\\
	&= \equivrep{(a_1 + a_2) + a_3, (b_1 + b_2) + b_3}{}\\
	&= \equivrep{a_1 + a_2, b_1 + b_2}{} + \equivrep{a_3, b_3}{}\\
	&= (\equivrep{a_1, b_1}{} + \equivrep{a_2, b_2}{}) + \equivrep{a_3, b_3}{}\\
	&= (i+j) + k
\end{align*}
dengan bebas memanfaatkan perilaku penjumlahan pada $\Nat$.
\end{proof}

Begitu pula, berikut cara membuktikan \emph{Invers aditif}:

\begin{proof}
Tetapkan $i \in \Int$, sehingga $i = \equivrep{a,b}{}$ untuk suatu $a,b \in
\Nat$. Misalkan $j = \equivrep{b,a}{} \in \Int$. Dengan memanfaatkan
perilaku bilangan asli, $(a+b) + 0 = 0 + (a+b)$, sehingga
$\tuple{a+b, b+a} \sim_\Int \tuple{0,0}$ menurut definisi, dan dengan demikian
$\equivrep{a+b, b+a}{} = \equivrep{0, 0}{} = 0_\Int$. Maka sekarang $i + j =
\equivrep{a,b}{}+\equivrep{b,a}{}=\equivrep{a+b, b+a}{}= \equivrep{0,
0}{} = 0_\Int$.
\end{proof}

Berikut pula pembuktian \emph{Distributivitas}:

\begin{proof}
Seperti di atas, tetapkan $i = \equivrep{a_1, b_1}{}$ dan $j =
\equivrep{a_2,b_2}{}$ dan $k = \equivrep{a_3, b_3}{}$. Sekarang:
\begin{align*}
	i \times (j + k) 
	&= \equivrep{a_1, b_1}{} \times (\equivrep{a_2,b_2}{} + \equivrep{a_3, b_3}{})\\
	&= \equivrep{a_1, b_1}{} \times \equivrep{a_2 + a_3,b_2+b_3}{}\\
	&= \equivrep{a_1  (a_2 + a_3) + b_1  (b_2+b_3), a_1  (b_2 + b_3) + b_1 (a_2 + a_3)}{}\\
	&= \equivrep{a_1 a_2 + a_1a_3 + b_1 b_2+b_1b_3, a_1 b_2 + a_1b_3 + a_2b_1 + a_3b_1}{}\\		
%		&= \equivrep{a_1 a_2 + b_1b_2 + a_1a_3 + b_1 b_2+b_1b_3, a_1 b_2 + a_1b_3 + a_2b_1 + a_3b_1}{}\\		
	&= \equivrep{a_1a_2 + b_1b_2, a_1b_2 + a_2b_1}{} + \equivrep{a_1a_3 + b_1b_3, a_1b_3 + a_3b_1}{}\\
	&= (\equivrep{a_1, b_1}{} \times \equivrep{a_2,b_2}{}) + (\equivrep{a_1, b_1}{} \times  \equivrep{a_3, b_3}{})\\
	&= (i \times j) + (i \times k)
\end{align*}
\end{proof}

Sebagai latihan, silakan buktikan kelima kondisi lainnya. Setelah
melakukannya, kita telah menunjukkan bahwa $\Int$ merupakan ring komutatif,
yakni bahwa penjumlahan dan perkalian (sebagaimana didefinisikan) berperilaku
sebagaimana semestinya.

\begin{prob}
Buktikan bahwa $\Int$ adalah ring komutatif.
\end{prob}

Namun, tugas kita belum selesai. Selain mendefinisikan penjumlahan dan
perkalian pada $\Int$, kita mendefinisikan suatu relasi urutan, $\leq$, dan
kita harus memeriksa bahwa relasi ini berperilaku sebagaimana semestinya.
Secara lebih rinci, kita harus menunjukkan bahwa $\Int$ merupakan ring
\emph{terurut}.\footnote{Ingat dari
	\olref[sfr][rel][ord]{def:linearorder} bahwa suatu urutan total
	adalah relasi yang refleksif, transitif, antisimetris, dan terhubung.
	Dalam konteks relasi urutan, keterhubungan kadang-kadang disebut
	\emph{trikotomi}, karena untuk setiap $a$ dan $b$ berlaku $a \leq b \lor a
	= b \lor a \geq b$.}

\begin{defn}
Suatu \emph{ring terurut} adalah ring komutatif yang juga dilengkapi
dengan suatu relasi urutan total, $\leq$, sedemikian sehingga:
\begin{align*}
	a \leq b &\lif a + c \leq b + c\\
	(a \leq b \land 0 \leq c) &\lif a \times c \leq b \times c
\end{align*}
\end{defn}

\begin{prob}
Buktikan bahwa $\Int$ adalah ring terurut.
\end{prob}

Seperti sebelumnya, menunjukkan bahwa~$\Int$, sebagaimana dibangun, merupakan
ring terurut adalah pekerjaan rutin tetapi panjang. Hal itu kita serahkan
kepada Anda.

Dengan demikian, bilangan bulat telah selesai ditangani. Namun, sekarang kita
perlu menunjukkan hal-hal yang sangat mirip mengenai bilangan rasional. Secara
khusus, sekarang kita perlu menunjukkan bahwa bilangan rasional membentuk
suatu \emph{lapangan} terurut berdasarkan definisi kita bagi $+$, $\times$,
dan $\leq$:
\begin{defn}\ollabel{orderedfield}
Suatu \emph{lapangan terurut} adalah ring terurut yang juga memenuhi:
\begin{align*}
	\emph{Invers perkalian}& & (\forall a \in S \setminus \{0\})(\exists b \in S) a\times b& = 1
\end{align*}
\end{defn}

Setelah Anda menunjukkan bahwa $\Int$ merupakan ring terurut, menunjukkan
bahwa $\Rat$ merupakan lapangan terurut adalah pekerjaan yang mudah tetapi
panjang.

\begin{prob}
Buktikan bahwa $\Rat$ adalah lapangan terurut.
\end{prob}

Setelah menangani bilangan bulat dan bilangan rasional, tinggal bilangan real
yang perlu ditangani. Secara khusus, kita perlu menunjukkan bahwa $\Real$
merupakan lapangan terurut yang \emph{lengkap}, yakni lapangan terurut yang
memiliki Sifat Kelengkapan. \olref[cuts]{realcompleteness} telah membuktikan
bahwa $\Real$ memiliki Sifat Kelengkapan. Namun, kita masih perlu menjalani
proses (yang melelahkan) untuk memeriksa bahwa $\Real$ merupakan lapangan
terurut.

Sebelum bergegas mengerjakan latihan panjang \emph{itu}, kita perlu memeriksa
beberapa hal yang lebih ``langsung''. Sebagai contoh, kita memerlukan jaminan
bahwa $\alpha + \beta$, sebagaimana didefinisikan, memang suatu
\emph{potongan}, untuk sembarang potongan $\alpha$ dan $\beta$. Berikut
pembuktian fakta tersebut:

\begin{proof}	
Karena $\alpha$ dan $\beta$ keduanya potongan, $\alpha + \beta = \Setabs{p
+ q}{p \in \alpha \land q \in \beta}$ merupakan himpunan bagian sejati
takkosong dari $\Rat$. Sekarang andaikan $x < p + q$ untuk suatu $p \in
\alpha$ dan $q \in \beta$. Maka $x - p < q$, sehingga $x - p \in \beta$,
dan $x = p + (x - p) \in \alpha + \beta$. Jadi $\alpha + \beta$ adalah
segmen awal dari $\Rat$. Terakhir, untuk setiap $p + q \in \alpha + \beta$,
karena $\alpha$ dan $\beta$ keduanya potongan, terdapat $p_1 \in \alpha$
dan $q_1 \in \beta$ sedemikian sehingga $p < p_1$ dan $q < q_1$; jadi $p +
q < p_1 + q_1 \in \alpha + \beta$; maka $\alpha + \beta$ tidak mempunyai
maksimum.
\end{proof}

Upaya serupa akan memungkinkan Anda memeriksa bahwa $\alpha - \beta$ dan
$\alpha \times \beta$ dan $\alpha \div \beta$ merupakan potongan (dalam kasus
terakhir, dengan mengabaikan kasus ketika $\beta$ merupakan potongan nol).
Namun, sekali lagi, hal itu kita serahkan kepada Anda.

\begin{prob}
Buktikan bahwa $\Real$ adalah lapangan terurut.
\end{prob}

Namun, ada satu rincian kecil yang perlu dirapikan. Dalam
\olref[cuts]{sec}, kita mengusulkan agar kita dapat menetapkan $\sqrt{2} =
\Setabs{p \in \Rat}{p < 0 \text{ atau }p^2 < 2}$. Akan tetapi, kita memang
perlu menunjukkan bahwa himpunan ini merupakan suatu \emph{potongan}. Berikut
pembuktian fakta tersebut:

\begin{proof}
Jelas bahwa himpunan ini merupakan segmen awal sejati takkosong dari bilangan
rasional; jadi cukup ditunjukkan bahwa himpunan ini tidak mempunyai maksimum.
Secara khusus, cukup ditunjukkan bahwa jika $p$ adalah bilangan rasional
positif dengan $p^2 < 2$ dan $q = \frac{2p+2}{p+2}$, maka $p < q$ dan $q^2 <
2$. Untuk melihat bahwa $p < q$, cukup perhatikan:
\begin{align*}
	p^2 &< 2\\
	p^2 + 2p &< 2 + 2p\\
	p(p + 2) &< 2 + 2p\\
	p &< \tfrac{2+2p}{p+2} = q
\end{align*}
Untuk melihat bahwa $q^2 < 2$, cukup perhatikan:
\begin{align*}
	p^2 &< 2\\
	2p^2 + 4p + 2 &< p^2 + 4p+ 4\\
	4p^2 + 8p + 4 &< 2(p^2 + 4p + 4)\\
	(2p+2)^2 & <2(p+2)^2\\
	\tfrac{(2p+2)^2}{(p+2)^2} &< 2\\
	q^2 &< 2
\end{align*}
\end{proof}
\end{document}
